\chapter{Geometria Geodésica e Estrutura do Radome}

A revisão 3 da arquitetura parte de um icosaedro regular proposto (parâmetros \texttt{GEO-001} e \texttt{GEO-002}), com $V_0=12$, $E_0=30$ e $F_0=20$. A candidata tetraédrica da ADR-012 usa uma subdivisão planar classe I de frequência 2: cada macroface equilátera é dividida pelos pontos médios das arestas em quatro faces receptoras equiláteras e coplanares, sem projetar os novos pontos para uma esfera. O envelope fechado verificado por \texttt{geometry/verify\_tetrahedral\_face\_geometry.py} permanece
\[
V=42,\qquad E=120,\qquad F=80,
\]
com consistência de Euler $V-E+F=2$ (\texttt{GEO-003}, \texttt{GEO-009} e \texttt{GEO-010}). Todas as faces receptoras externas passam a ter três arestas de \SI{2.0000}{m} (\texttt{GEO-004} e \texttt{GEO-015}), enquanto cada macroaresta mede \SI{4.0000}{m}. O envelope facetado possui raio circunscrito de \SI{3.8042}{m} e inraio de \SI{3.0230}{m} (\texttt{GEO-017}). Esses valores substituem as dimensões da esfera projetada para esta candidata.

Um tetraedro regular literal, com as seis arestas de \SI{2}{m}, não pode ser repetido nesse conjunto fechado voltado para dentro: o verificador detecta 120 interpenetrações volumétricas (\texttt{GEO-019}). A interpretação adotada mantém apenas as três arestas externas com \SI{2}{m} e usa um ápice interno independente a \SI{0.75}{m} do centro de cada face (\texttt{GEO-021}). As 80 células tetraédricas rasas não apresentam colisão volumétrica (\texttt{GEO-020}) e suas arestas laterais medem \SI{1.3769}{m}. Os ápices vizinhos ficam separados por \SI{0.5435}{m} a \SI{1.1547}{m}; as duas paredes Faraday independentes que partem de cada aresta externa compartilhada delimitam, portanto, um corredor intercelular para rotas separadas de energia, fibra óptica, equipotencialização e manutenção (\texttt{MEC-003} e \texttt{MEC-005}). A construção reserva um núcleo interno livre com raio mínimo de \SI{2.2730}{m} (\texttt{GEO-022}). O corte circular e o anel de apoio anteriores pertencem à candidata projetada e são históricos; o C2 permanece aberto até verificar uma borda formada por módulos completos e uma nova transição civil.

O esqueleto mecânico, os painéis dielétricos e os módulos RF internos possuem geometrias diferentes, porém coordenadas. Uma face é um setor mecânico e computacional, não uma região eletromagneticamente isolada. A sobreposição angular entre faces vizinhas é necessária para calibração, interpolação e estimação de direção.

A abertura RF externa ocupa a face triangular não blindada. Para qualquer faixa declarada como dual-pol, dois componentes ortogonais e simultâneos seguem a base tangente local $\mathbf{t}_u,\mathbf{t}_v$ e usam alimentações e invólucros ADC/ASIC blindados independentes; a coerência e os produtos polarimétricos continuam sujeitos à calibração na mesma faixa exigida pelo C1. As outras três faces tetraédricas são paredes condutivas da gaiola de Faraday. As passagens usam feedthroughs filtrados e os invólucros se ligam a uma placa fria condutiva sem tornar a abertura dielétrica o caminho estrutural primário.


As junções internas formam caminhos contínuos de carga entre os membros da estrutura triangular e as faces vizinhas. Costuras estruturais vedadas, reforços locais e um anel ou placa nodal interna transferem cargas de vento, vibração e manutenção sem depender apenas da pele dielétrica. O conjunto deve ser verificado por análise estrutural e ensaios ambientais.


O modelo Blender torna explícita a convenção de montagem (\texttt{MEC-006} e \texttt{MEC-007}): a face RF triangular está na superfície externa do radome, enquanto o volume da pirâmide e seu vértice apontam para dentro. Uma haste parte do ápice interno, atravessa o centro da face e prossegue para fora como eixo comum dos booms, colinear com a normal externa dessa face. Os elementos VHF e UHF são transversais à normal, mutuamente ortogonais, e a montagem cruzada completa é girada em 45 graus ao redor da normal em relação à base tangente local. Essa rotação mecânica não transforma os canais de faixas diferentes em um par polarimétrico coerente.




\section{Solução estrutural}
A carga da Yagi externa é transferida pela haste de base para a moldura triangular e, em seguida, para o nó interno e as junções das faces vizinhas. A pele dielétrica é um fechamento RF e ambiental, não o único caminho estrutural primário. O modelo estrutural deve incluir pressão de vento na Yagi e no radome, flexão da haste, massa dos equipamentos, expansão térmica, vibração, cargas de manutenção e caminhos de corrente de descargas. Isso fecha a lacuna arquitetural entre um conceito de antena e um módulo de face sustentado, embora ainda sejam necessários análise por elementos finitos, ensaio modal e qualificação climática.

\begin{figure}[H]
\centering
\includegraphics[width=0.98\textwidth]{pt-BR/fig16_tetrahedral_face_cluster.png}
\caption{Sete células tetraédricas rasas contíguas. Os membros laranja e ciano mostram as direções VHF/UHF mutuamente ortogonais a 45° e 135° em cada base tangente local; as setas verdes mostram as normais externas das faces, que definem as direções dos booms. O envelope verde aramado é o núcleo interno de serviço reservado.}
\label{fig:fig16_tetrahedral_face_cluster}
\end{figure}

\begin{figure}[H]
\centering
\includegraphics[width=0.98\textwidth]{pt-BR/fig11_face_yagi_camadas.png}
\caption{Montagem de face. O boom escuro segue a normal externa desde o ápice interno e atravessa o centro da abertura triangular de dois metros; os elementos laranja e azul são mutuamente ortogonais e girados conjuntamente em 45° ao redor dessa normal. As placas coloridas à direita representam a pilha de materiais e eletrônica.}
\label{fig:fig11_face_yagi_camadas}
\end{figure}

\begin{figure}[H]
\centering
\includegraphics[width=0.95\textwidth]{pt-BR/fig02_zonamento_radome.png}
\caption{Zoneamento funcional: os setores coloridos representam as famílias híbridas de faixas e o centro branco representa sincronismo, calibração, processamento e serviços compartilhados na face. Os limites dos setores são alocações arquiteturais, não fronteiras eletromagnéticas rígidas.}
\label{fig:fig02_zonamento_radome}
\end{figure}

\begin{figure}[H]
\centering
\includegraphics[width=0.98\textwidth]{pt-BR/fig13_radome_blender.png}
\caption{Vista externa no Blender do radome e da face explodida. O boom combinado VHF/UHF é colinear com a normal da face; os elementos transversais são ortogonais e girados em 45° ao redor dessa normal.}
\label{fig:fig13_radome_blender}
\end{figure}

\begin{figure}[H]
\centering
\includegraphics[width=0.98\textwidth]{pt-BR/fig14_radome_interior_blender.png}
\caption{Vista interna de inspeção no Blender com casca transparente, caminhos estruturais de carga, troncos de serviço e módulos blindados representativos.}
\label{fig:fig14_radome_interior_blender}
\end{figure}
